% =========================================================
% Notes: Number Systems as L-Structures
% File: volume-vi/book-algebra/algebraic-structures/notes/notes-number-systems-as-structures.tex
%
% Purpose: Reframes N, Z, Q, R as concrete models of progressively
% richer signatures. Links the concrete axiom table in Volume II to
% the abstract blueprint hierarchy in this chapter.
% =========================================================

\subsection{\texorpdfstring{The Number Systems as $\mathcal{L}$-Structures}{The Number Systems as L-Structures}}
\label{sec:number-systems-as-structures}

% ---------------------------------------------------------
\subsection*{The Arithmetic Signature}

The number systems $\mathbb{N}, \mathbb{Z}, \mathbb{Q}, \mathbb{R}$
all live in the same signature --- they differ in which axiom
\emph{sentences} they satisfy, not in which symbols they carry.

\begin{axiombox}{The Arithmetic Signature $\mathcal{L}_{\mathrm{arith}}$}

\[
\mathcal{L}_{\mathrm{arith}}
\;=\;
\bigl\{\;
  \underbrace{+,\;\cdot}_{\text{arity 2}},\quad
  \underbrace{-,\;(\cdot)^{-1}}_{\text{arity 1}},\quad
  \underbrace{0,\;1}_{\text{constants}},\quad
  \underbrace{{\leq}}_{\text{relation, arity 2}}
\;\bigr\}.
\]

\medskip
\begin{tabular}{@{} l l l @{}}
\toprule
\textbf{Symbol} & \textbf{Arity} & \textbf{Intended interpretation} \\
\midrule
$+$        & 2 (function) & addition \\
$\cdot$    & 2 (function) & multiplication \\
$-$        & 1 (function) & additive inverse / negation \\
$(\cdot)^{-1}$ & 1 (function) & multiplicative inverse \\
$0$        & 0 (constant) & additive identity \\
$1$        & 0 (constant) & multiplicative identity \\
${\leq}$   & 2 (relation) & order \\
\bottomrule
\end{tabular}
\end{axiombox}

\begin{remark}
Every number system is an $\mathcal{L}_{\mathrm{arith}}$-structure:
it carries all these symbols.
What distinguishes $\mathbb{N}$ from $\mathbb{Z}$ from $\mathbb{R}$
is not the signature but the \emph{theory} --- the set of first-order
sentences in $\mathcal{L}_{\mathrm{arith}}$ that hold in each model.

In particular, $(\cdot)^{-1}$ is a symbol in the signature of
$\mathbb{Z}$, but the sentence
\[
\forall x\;(x \neq 0 \Rightarrow x \cdot x^{-1} = 1)
\]
is \emph{false} in $\mathbb{Z}$.
The symbol is interpreted as a total function on $\mathbb{Z}$
(defined arbitrarily on $0$ and on nonzero elements where the true
inverse does not exist), but the field axiom it satisfies in
$\mathbb{Q}$ and $\mathbb{R}$ fails.
This is the model-theoretic way of saying: $\mathbb{Z}$ is not a field.
\end{remark}

% ---------------------------------------------------------
\subsection*{Each Number System as a Model}

\begin{tcolorbox}[colback=gray!6, colframe=gray!40, arc=2pt,
  left=6pt, right=6pt, top=4pt, bottom=4pt,
  title={\small\textbf{Number Systems as Models of
         $\mathcal{L}_{\mathrm{arith}}$}},
  fonttitle=\small\bfseries]
\small
\renewcommand{\arraystretch}{1.35}
\begin{tabular}{@{} l l l @{}}
\toprule
\textbf{Structure} & \textbf{Carrier} & \textbf{Theory (axiom sentences satisfied)} \\
\midrule
$\mathcal{N} = (\mathbb{N}, +, \cdot, -, (\cdot)^{-1}, 0, 1, \leq)$
  & $\mathbb{N}$
  & Semiring axioms + well-ordering.
    $-$ and $(\cdot)^{-1}$ defined but inverse axioms fail. \\[4pt]
$\mathcal{Z} = (\mathbb{Z}, +, \cdot, -, (\cdot)^{-1}, 0, 1, \leq)$
  & $\mathbb{Z}$
  & Commutative ordered ring axioms.
    Additive inverse axiom holds; multiplicative inverse axiom fails. \\[4pt]
$\mathcal{Q} = (\mathbb{Q}, +, \cdot, -, (\cdot)^{-1}, 0, 1, \leq)$
  & $\mathbb{Q}$
  & Ordered field axioms.
    Both inverse axioms hold; least upper bound axiom fails. \\[4pt]
$\mathcal{R} = (\mathbb{R}, +, \cdot, -, (\cdot)^{-1}, 0, 1, \leq)$
  & $\mathbb{R}$
  & Complete ordered field axioms.
    All axioms hold; $\mathbb{R}$ is the unique (up to isomorphism)
    complete ordered field. \\
\bottomrule
\end{tabular}
\end{tcolorbox}

\begin{remark}[The abstract algebra framing of the extension hierarchy]
From the blueprint hierarchy in this chapter:
\[
\underbrace{\mathbb{N}}_{\text{semiring}}
\;\subset\;
\underbrace{\mathbb{Z}}_{\text{comm.\ ring}}
\;\subset\;
\underbrace{\mathbb{Q}}_{\text{ordered field}}
\;\subset\;
\underbrace{\mathbb{R}}_{\text{complete ordered field}}.
\]
Each inclusion is a \emph{substructure embedding}: the smaller system
embeds into the larger as an $\mathcal{L}_{\mathrm{arith}}$-substructure
(operations and constants are preserved, but more sentences become true
in the larger model).

More precisely:
$\mathbb{N} \hookrightarrow \mathbb{Z}$ adds the solution $-n$ to every
equation $n + x = 0$.
$\mathbb{Z} \hookrightarrow \mathbb{Q}$ adds the solution $p/q$ to every
equation $qx = p$ with $q \neq 0$.
$\mathbb{Q} \hookrightarrow \mathbb{R}$ adds a limit point to every
bounded increasing sequence --- it is a \emph{completion}, not a purely
algebraic extension.
\end{remark}

% ---------------------------------------------------------
\subsection*{The Blueprint Position of Each Number System}

\begin{tcolorbox}[colback=gray!6, colframe=gray!40, arc=2pt,
  left=6pt, right=6pt, top=4pt, bottom=4pt,
  title={\small\textbf{Number Systems in the Blueprint Hierarchy}},
  fonttitle=\small\bfseries]
\small
\renewcommand{\arraystretch}{1.3}
\begin{tabular}{@{} l l l @{}}
\toprule
\textbf{System} & \textbf{Blueprint level} & \textbf{What it adds} \\
\midrule
$\mathbb{N}$
  & Commutative semiring
  & Two operations, distributivity, no additive inverses. \\
$\mathbb{Z}$
  & Commutative ordered ring
  & Additive inverses (Blueprint 8 = Ring, commutative). \\
$\mathbb{Q}$
  & Ordered field
  & Multiplicative inverses (Blueprint 9 = Field). \\
$\mathbb{R}$
  & Complete ordered field
  & Least upper bound property (analytic, not purely algebraic). \\
\bottomrule
\end{tabular}
\end{tcolorbox}

\begin{remark}[Where completeness lives]
The least upper bound property is \emph{not} expressible as a single
first-order sentence in $\mathcal{L}_{\mathrm{arith}}$ --- it is a
second-order axiom schema (quantifying over subsets).
This is why $\mathbb{R}$ cannot be characterised up to isomorphism
by any first-order theory: by the compactness theorem, any first-order
theory satisfied by $\mathbb{R}$ also has non-standard models
(the hyperreals).
The categorical characterisation --- $\mathbb{R}$ is the unique complete
ordered field up to isomorphism --- requires the full second-order
least upper bound axiom.
This is the boundary between algebra and analysis.
\end{remark}

\begin{remark}[Derived properties are sentences, not axioms]
Properties like $a \cdot 0 = 0$, $(-a)b = -(ab)$, and
$ab = 0 \Rightarrow a = 0 \lor b = 0$ (no zero divisors in $\mathbb{Z}$)
are \emph{theorems}: first-order sentences derivable from the ring axioms.
They hold in every model of those axioms, not just in $\mathbb{Z}$.
From the model-theoretic perspective, the ring axioms are a theory $T$,
and these properties are elements of the \emph{deductive closure}
$\mathrm{Th}(T)$ --- they are true in every model of $T$.
See Volume II \S\ref{sec:number-systems-axioms} for the concrete
derivations.
\end{remark}
